\chapter{MOTOR DE OTIMIZAÇÃO NO BACKEND}

Foi integrado ao compilador um Motor de Otimização Independente de Máquina, atuando no backend diretamente sobre a Árvore de Sintaxe Abstrata (AST) antes da geração de código.

O módulo está implementado em \texttt{src/optimizer.py} e empacotado no componente \texttt{C\_Otimizador} (\texttt{src/componente\_otimizador.py}), obedecendo à interface \texttt{IOtimizador}.

\section{Técnicas de Otimização Implementadas}

A classe \texttt{ASTOptimizer} implementa transformações semânticas que preservam o comportamento observável do programa, reduzindo a profundidade da árvore sintática, eliminando nós supérfluos e computando valores constantes em tempo de compilação:

\begin{enumerate}
    \item Dobra de Constantes (\textit{Constant Folding}):
    expressões aritméticas, lógicas e relacionais formadas exclusivamente por literais estáticos são avaliadas previamente pelo compilador, substituindo subárvores inteiras por um nó literal único:
    \begin{itemize}
        \item $2 + 3 \times 4 \;\Longrightarrow\; 14$
        \item $\texttt{"Compiladores "} + \texttt{"2026.2"} \;\Longrightarrow\; \texttt{"Compiladores 2026.2"}$
        \item $\texttt{true } \&\& \texttt{ false} \;\Longrightarrow\; \texttt{false}$
    \end{itemize}

    \item Eliminação de Código Morto (\textit{Dead Code Elimination}):
    comandos e blocos inalcançáveis em tempo de execução são podados da AST:
    \begin{itemize}
        \item Blocos \texttt{if (false) \{ ... \} else \{ B \}} têm o teste condicional e o ramo principal descartados, retendo-se diretamente o ramo \texttt{B}.
        \item Blocos \texttt{if (true) \{ A \} else \{ ... \}} são simplificados diretamente para o corpo \texttt{A}.
        \item Enunciados subsequentes a instruções incondicionais de desvio (\texttt{return}, \texttt{break}, \texttt{continue}) no mesmo bloco sequencial são eliminados.
    \end{itemize}

    \item Simplificações Algébricas (\textit{Algebraic Identities}):
    aplica propriedades matemáticas para simplificar termos envolvendo identidades:
    \begin{itemize}
        \item $x + 0 \;\Longrightarrow\; x \quad\text{e}\quad 0 + x \;\Longrightarrow\; x$
        \item $x - 0 \;\Longrightarrow\; x$
        \item $x \times 1 \;\Longrightarrow\; x \quad\text{e}\quad 1 \times x \;\Longrightarrow\; x$
        \item $x \times 0 \;\Longrightarrow\; 0 \quad (\text{se } x \text{ for puro})$
        \item $x / 1 \;\Longrightarrow\; x$
    \end{itemize}
\end{enumerate}

\section{Garantia de Pureza e Ponto Fixo}

Para garantir que a eliminação de expressões algébricas não descarte acidentalmente efeitos colaterais relevantes (como chamadas de função com I/O ou manipulação de variáveis globais), o otimizador inclui a verificação estática de pureza através do método \texttt{is\_pure(node)}. Uma subárvore só é considerada pura se contiver estritamente variáveis, constantes ou operadores sem efeitos secundários.

As passadas de otimização são executadas em um laço de ponto fixo (\textit{fixed-point iteration}): a árvore continua sendo reprocessada recursivamente enquanto transformações forem geradas, até que o número de nós se estabilize ou o limite de segurança de iterações seja atingido.

\section{Métricas e Ativação no Compilador}

A otimização pode ser habilitada opcionalmente:
\begin{itemize}
    \item Linha de Comando (CLI): através da flag \texttt{-O} ou \texttt{--optimize}:
    \begin{lstlisting}[language=bash]
python main.py tests/oficiais/teste02_fatorial_fibonacci_par.minipar -O
    \end{lstlisting}
    \item Interface Web Didática: através do seletor \textit{"Otimizar AST"}, que exibe um relatório com o número de nós antes e depois da transformação, o percentual de redução obtido e o registro textual de todas as regras aplicadas.
\end{itemize}
